\section{Why Speed Matters, and Why Speed Is Not the Same as Haste}

\subsection{The Velocity of Uncertainty Reduction}

Startups are often associated with urgency. That association is justified: in some markets, the window for assembling talent, reaching customers, and establishing a position is brief. Knowledge also has a half-life. A team that can test 10 plausible ideas while another tests 1 will learn much faster.

But the productive variable is not the number of hours spent frantically moving. It is the rate at which valuable uncertainty is removed and useful capability is accumulated.

Suppose a founder spends a month writing a 100-page strategy document without speaking to a customer. Another spends the month selling a simple version of the proposed service, studying the rejection reasons, and improving the offer. The 2nd founder may be learning at a higher rate, even if their documents and website look less impressive.

Now consider a medical device or an academic assessment used to certify competence. An error can harm patients or misrepresent someone's abilities. Speed without verification is not entrepreneurial excellence. The faster team is the one that finds reliable ways to evaluate safety and validity, not the one that simply ships untested claims sooner.

The advantage of a small organization is often its capacity to keep decision cycles short. The founder can notice a problem, investigate it, change the approach, and learn from the result without waiting for committees. As the business grows, this advantage can disappear unless decision rights, information flows, and incentives are deliberately designed.

\textbf{The important lesson for day one is simple: be fast in learning, explicit about what must be verified, and slow to make irreversible commitments without evidence.}
